
\subsubsection{2.1 Parameter-Efficient Fine-Tuning and Adaptive Rank}

LoRA \citep{hu2021lora} constrains weight updates to low-rank matrices, reducing trainable parameters by orders of magnitude. The original formulation assigns a uniform rank r across all target layers, which is acknowledged as potentially suboptimal.

\textbf{AdaLoRA} \citep{zhang2023adalora} extends LoRA with singular value decomposition of weight update matrices. Singular values are adaptively pruned during training via an importance score, achieving better parameter efficiency than uniform rank. The rank allocation is a product of the training run and must be re-derived per task.

\textbf{DyLoRA} \citep{valipour2022dylora} addresses the rank search problem by training across multiple rank values simultaneously via stochastic rank sampling. Like AdaLoRA, DyLoRA produces the rank distribution as a result of training rather than providing it as an input.

\textbf{Act-LoRA} (unverified citation; see notes) uses layer-wise L2 activation norms for binary layer selection --- whether to include a layer in LoRA at all --- rather than continuous rank magnitude allocation, and does not characterize the stability or threshold invariance of the activation signal.

The present work differs from all three in that the proposed signal is measured before training, from a single forward pass, with no gradient computation.

\subsubsection{2.2 Activation Sparsity in Large Language Models}

Li et al. (2023) characterize the Lazy Neuron Phenomenon: neurons become increasingly sparse during training, with feed-forward network (FFN) activations exhibiting concentrated activity. Mirzadeh et al. (2024) demonstrate in ''ReLU Strikes Back'' that modern LLMs achieve substantial sparsity at inference time. Liu et al. (2025) provide the TEAL framework, a training-free approach to exploiting activation sparsity via magnitude-based thresholding, reporting 40--50\% model-wide sparsity on LLaMA-2/3-8B with minimal degradation.

Szatkowski et al. (2025) identify universal structural properties across LLM families, finding that sparsity patterns exhibit cross-model consistency. This prior work motivates the cross-distribution stability analysis in the present study; the ICC(3,k) = 0.9846 result obtained here strongly confirms distribution-invariance for LLaMA-3.1-8B specifically.

The existing literature characterizes the existence and approximate magnitude of activation sparsity in large transformers. It does not establish whether the layer-wise sparsity fingerprint is stable enough --- across calibration distributions and threshold choices --- to serve as a pre-training prior for rank allocation decisions. That characterization is the central empirical contribution of this work.

\subsubsection{2.3 Intrinsic Dimensionality and Layer Structure}

Aghajanyan et al. (2021) demonstrate that pre-trained language models have low intrinsic dimension for fine-tuning. Their structure-aware projection method (SAID) significantly outperforms dimension-insensitive alternatives, providing a theoretical bridge: if activation sparsity proxies intrinsic dimension, it could directly inform rank allocation. The present work does not test this bridge empirically; H-M3 is designed to do so.

Clark et al. (2019) and Tenney et al. (2019) demonstrate a linguistic hierarchy in transformer layers, with syntactic processing concentrated in early layers and semantic integration in later ones. The depth gradient observed here --- early layers most sparse, deep layers least sparse --- is consistent with this specialization, though the mechanistic link is not established by this work.

\begin{table*}[htbp]
\centering
\resizebox{\dimexpr 2\columnwidth+\columnsep\relax}{!}{%
\begin{tabular}{lllll}
\toprule
Method & Signal & Measurement Cost & Allocation Type & Stability Characterized \\
\midrule
AdaLoRA (Zhang et al., 2023) & SVD importance & Full training run & Continuous & No \\
DyLoRA (Valipour et al., 2023) & Rank sampling & Full training run & Continuous & No \\
TEAL (Liu et al., 2025) & Activation magnitude & Inference pass & Pruning & No \\
This work & Layer sparsity & Single forward pass & Prior (continuous) & Yes \\
\bottomrule
\end{tabular}
}
\end{table*}

